A name written after a dot is an \emph{attribute reference},
and one rule fixes what it denotes \citep[Section 9.3]{PSF26tut}.
We state it for classes that may be built on other classes,
the subject of Section \ref{sec.inheritance}.

\begin{defi}\label{def.attributeReference}
 Every class $C$ has a \emph{linearisation},
 which the language calls its method resolution order:
 the tuple \codehl{C.\_\_mro\_\_} of classes $(C_0, \dots, C_n)$,
 where $C_0 = C$ and $C_n$ is the class \codehl{object}.
 \begin{enumerate}[label={\emph{\roman{*}}.}]
  \item The reference \codehl{C.name} denotes the attribute \codehl{name}
        of the first of $C_0, \dots, C_n$ that has one.
  \item For an instance $x$ of $C$,
        the reference \codehl{x.name} denotes the attribute \codehl{name} of $x$,
        if $x$ has one.
        Otherwise it denotes \codehl{C.name};
        and when \codehl{C.name} is a function $f$,
        it denotes the function $u \mapsto f(x, u)$ of the remaining arguments $u$,
        which is called a \emph{method} of $x$.
 \end{enumerate}
 A reference that finds no attribute raises \codehl{AttributeError}.
\end{defi}

\codehl{AggregateBook} is built on no other class,
and its linearisation is itself followed by \codehl{object}.
By part \emph{ii.}, the call \codehl{x.name(u)} of a method
is the call \codehl{C.name(x, u)} of a function of the class:
the parameter \codehl{self} of Listings \ref{listing.aggregateStorage}
and \ref{listing.aggregateInit} receives the object written to the left of the dot.
The functions are thus held once, by the class,
and an instance holds only what distinguishes it from the others.

An assignment is not a search.
A reference looks in the instance and then in the classes of Definition \ref{def.attributeReference},
whereas \codehl{x.name = value} binds the name in the instance, whatever the classes hold,
and \codehl{del x.name} removes it from the instance.

The operations of the language are attribute references too,
made on the type of their operand \citep[Section 3.3]{PSF26lan}:
\codehl{len(x)} is the call \codehl{type(x).\_\_len\_\_(x)}.
A name with two underscores on either side is reserved for this use,
and a class takes part in an operation by binding the function of that name,
a \emph{special method}.
Table \ref{tab.operationsAsCalls} lists the operations the package relies on.
A class that binds none of them finds those of \codehl{object},
which prints an instance as the name of its class and an address,
and compares two instances by identity.

\begin{table}[h!]
 \centering
 \begin{tabular}{@{}lll@{}}
  \toprule
  expression & is the call & in the package \\
  \midrule
  \codehl{C(u)} & \codehl{C.\_\_init\_\_(x, u)}, on a new $x$ & every class \\
  \codehl{repr(x)} & \codehl{type(x).\_\_repr\_\_(x)} & written in \codehl{AggregateBook} \\
  \codehl{len(x)} & \codehl{type(x).\_\_len\_\_(x)} & written in \codehl{EventJournal} \\
  \codehl{x == y} & \codehl{type(x).\_\_eq\_\_(x, y)} & generated, Section \ref{sec.dataClasses} \\
  \codehl{hash(x)} & \codehl{type(x).\_\_hash\_\_(x)} & generated, Section \ref{sec.dataClasses} \\
  \bottomrule
 \end{tabular}
 \caption{Operations of the language as calls of special methods.}
 \label{tab.operationsAsCalls}
\end{table}

Listing \ref{listing.methodIsFunction} checks Definition \ref{def.attributeReference}
on \codehl{best\_price}, the function of Listing \ref{listing.aggregateStorage}:
the name is an attribute of the class and not of the book,
the two ways of calling it return the same price,
and the method carries the book and the function of the class.

\begin{lstlisting}[caption={A method is a function of the class, with the instance prepended.}, label=listing.methodIsFunction]
from unito26.lob.messages import BUY
from unito26.lob.orderbook import AggregateBook

book = AggregateBook()
book.set_size(BUY, 99, 100)

assert AggregateBook.__mro__ == (AggregateBook, object)
assert "best_price" not in vars(book)
assert "best_price" in vars(AggregateBook)
assert book.best_price(BUY) == AggregateBook.best_price(book, BUY) == 99

method = book.best_price
assert method.__self__ is book
assert method.__func__ is AggregateBook.best_price
assert repr(book) == type(book).__repr__(book)
\end{lstlisting}
